\documentclass{article}
\usepackage{geometry}
\geometry{margin=1.5cm, vmargin={0pt,1cm}}
\setlength{\topmargin}{-1cm}
\setlength{\headheight}{12.5pt}
\setlength{\paperheight}{29.7cm}
\setlength{\textheight}{25.3cm}

% useful packages.
% \usepackage[UTF8]{ctex}
\usepackage{fancyhdr}
\usepackage{layout}

% import common macros
% % % % % % % % % % % % % % % % % % % % % % % % % % % % % % % %
% Common Macros for LaTeX
% % % % % % % % % % % % % % % % % % % % % % % % % % % % % % % %

% Useful packages
\usepackage{amsfonts}
\usepackage{amsmath}
\usepackage{amssymb}
\usepackage{amsthm}
\usepackage{enumerate}
\usepackage{graphicx}
\usepackage{multicol}
\usepackage[ruled,linesnumbered]{algorithm2e}

% Math operators and common symbols
\newcommand{\dif}{\mathrm{d}}
\newcommand{\innerProd}[1]{\left\langle #1 \right\rangle}
\newcommand{\difFrac}[2]{\frac{\dif #1}{\dif #2}}
\newcommand{\pdfFrac}[2]{\frac{\partial #1}{\partial #2}}
\newcommand{\T}{\mathrm{T}}
\newcommand{\norm}[1]{\left\| #1 \right\|}
\newcommand{\abs}[1]{\left| #1 \right|}

% Vector and Matrix shortcuts
\newcommand{\vecTwo}[2]{
  \begin{bmatrix}
    #1 \\
    #2
\end{bmatrix}}
\newcommand{\vecThree}[3]{
  \begin{bmatrix}
    #1 \\
    #2 \\
    #3
\end{bmatrix}}
\newcommand{\vecTwoH}[2]{
  \begin{bmatrix}
    #1 & #2
  \end{bmatrix}
}
\newcommand{\vecThreeH}[3]{
  \begin{bmatrix}
    #1 & #2 & #3
  \end{bmatrix}
}

% Common fields
\newcommand{\R}{\mathbb{R}}
\newcommand{\C}{\mathbb{C}}
\newcommand{\Z}{\mathbb{Z}}
\newcommand{\N}{\mathbb{N}}

% Solutions and proofs
\newenvironment{solution}{
\begin{proof}[解]}{
\end{proof}}

% Utility to get environment variables (requires lualatex)
\newcommand{\getEnv}[2]{\directlua{tex.print(os.getenv("#1") or "#2")}}


% Name and uID
\newcommand{\name}{\getEnv{NAME}{NAME}}
\newcommand{\uid}{\getEnv{UID}{UID}}
\newcommand{\course}{Point Set Topology}
\newcommand{\hwNum}{4}

\begin{document}

\pagestyle{fancy}
\fancyhead{}
\lhead{\zhstr{\name} (\uid)}
\chead{\course\ \#\hwNum}
\rhead{\today}

\section{Question 1}

Let
\[
A=\{(x,\sin(1/x)):x\in(0,1]\}\subset\R^2
\]
with the standard topology on $\R^2$. Compute $A^{\circ}$, $\overline{A}$, $LP(A)$, and $\overline{A}\cap(\R^2\setminus A)$.

\begin{solution}
The set $A$ is the graph of a function over $(0,1]$, so it has empty interior in $\R^2$:
\[
A^{\circ}=\varnothing.
\]

For closure, note first that $A\subset\overline A$. For $x\to0^+$, the values $\sin(1/x)$ oscillate in $[-1,1]$.
Given any $y\in[-1,1]$, choose $t_n\to\infty$ with $\sin t_n\to y$ and put $x_n=1/t_n$. Then $x_n\to0^+$ and
\[
(x_n,\sin(1/x_n))=(x_n,\sin t_n)\to(0,y).
\]
So $\{0\}\times[-1,1]\subset\overline A$.
Conversely, if a sequence in $A$ converges to $(a,b)$, then either $a>0$ and continuity of $x\mapsto\sin(1/x)$ on $(0,1]$ gives $b=\sin(1/a)$, or $a=0$ and necessarily $b\in[-1,1]$. Hence
\[
\overline A=A\cup(\{0\}\times[-1,1]).
\]

Now compute limit points.
Every point of $A$ is a limit point of $A$ because near each $x_0\in(0,1]$ there are infinitely many nearby $x\neq x_0$ in $(0,1]$, hence nearby points of the graph.
Every point of $\{0\}\times[-1,1]$ is also a limit point by the construction above.
There are no other limit points, so
\[
LP(A)=A\cup(\{0\}\times[-1,1]).
\]

Therefore
\[
\overline A\cap(\R^2\setminus A)=\{0\}\times[-1,1].
\]
\end{solution}

\section{Question 2}

Let $X$ be a set and $f:\mathcal P(X)\to\mathcal P(X)$ satisfy:
\[
\text{(i) }f(\varnothing)=\varnothing,\quad
\text{(ii) }A\subseteq f(A),\quad
\text{(iii) }f(f(A))=f(A),\quad
\text{(iv) }f(A\cup B)=f(A)\cup f(B).
\]
Define
\[
\mathcal{T}_f:=\{X\setminus C:\ C\subseteq X,\ f(C)=C\}.
\]
Prove: (1) $\mathcal{T}_f$ is a topology on $X$; (2) closure in $\mathcal{T}_f$ is exactly $f(A)$.

\begin{solution}
\textbf{Step 1: closed-set viewpoint.}
Let
\[
\mathcal C_f:=\{C\subseteq X:f(C)=C\}.
\]
Then $\mathcal{T}_f$ is exactly the family of complements of sets in $\mathcal C_f$.
So it is enough to show $\mathcal C_f$ is a closed-set system.

We first prove monotonicity of $f$.
If $A\subseteq B$, then $B=A\cup B$, so
\[
f(B)=f(A\cup B)=f(A)\cup f(B),
\]
hence $f(A)\subseteq f(B)$.

Now verify closed-set axioms:
\[
f(X)=f(X\cup\varnothing)=f(X)\cup f(\varnothing)=f(X),
\]
trivial; and by (ii), $X\subseteq f(X)\subseteq X$, so $f(X)=X$. Thus $X\in\mathcal C_f$.
Also $\varnothing\in\mathcal C_f$ by (i).

If $C_1,C_2\in\mathcal C_f$, then
\[
f(C_1\cup C_2)=f(C_1)\cup f(C_2)=C_1\cup C_2,
\]
so finite unions stay in $\mathcal C_f$.

If $\{C_i\}_{i\in I}\subseteq\mathcal C_f$ and $C=\bigcap_{i\in I}C_i$, then $C\subseteq C_i$ for all $i$, hence by monotonicity
\[
f(C)\subseteq f(C_i)=C_i\quad(\forall i),
\]
so $f(C)\subseteq\bigcap_i C_i=C$. By (ii), $C\subseteq f(C)$. Therefore $f(C)=C$ and $C\in\mathcal C_f$.

So $\mathcal C_f$ is exactly the family of closed sets of a topology, and therefore $\mathcal{T}_f$ is a topology.

\textbf{Step 2: closure is $f(A)$.}
Take any $A\subseteq X$.
By (iii),
\[
f(f(A))=f(A),
\]
so $f(A)$ is $\mathcal{T}_f$-closed; by (ii), $A\subseteq f(A)$.
Hence $f(A)$ is a closed superset of $A$.

Now let $C$ be any $\mathcal{T}_f$-closed set with $A\subseteq C$.
Then $f(C)=C$. By monotonicity,
\[
f(A)\subseteq f(C)=C.
\]
So $f(A)$ is contained in every closed superset of $A$, i.e. it is the closure:
\[
\overline A^{\,\mathcal{T}_f}=f(A).
\]
\end{solution}

\section{Question 3}

Consider
\[
f:[0,1)\to S^1,\qquad f(t)=(\cos(2\pi t),\sin(2\pi t)).
\]
Show $f$ is not a homeomorphism (with standard subspace topologies). Is there any homeomorphism between $[0,1)$ and $S^1$?

\begin{solution}
The map $f$ is continuous and bijective.
If it were a homeomorphism, then $[0,1)$ and $S^1$ would be homeomorphic.
But compactness is a topological invariant:
\begin{itemize}
\item $S^1$ is compact (closed and bounded in $\R^2$).
\item $[0,1)$ is not compact in $\R$.
\end{itemize}
Contradiction. So $f$ is not a homeomorphism.

Consequently, there is no homeomorphism between $[0,1)$ and $S^1$ at all.
\end{solution}

\section{Question 4}

Fix a topological space $(X,\mathcal{T})$. A set $A\subseteq X$ is regularly open if
\[
A=(\overline A)^\circ.
\]
(i) Give an open set in $\R$ (standard topology) that is not regularly open, and characterize regularly open sets in $\R$.
(ii) Prove that $(\overline A)^\circ$ is regularly open for any $A\subseteq X$.

\begin{solution}
\textbf{(i) Example and characterization in $\R$.}

Take
\[
U=(0,1)\cup(1,2).
\]
Then $U$ is open, but
\[
\overline U=[0,2],\qquad (\overline U)^\circ=(0,2)\neq U.
\]
So $U$ is not regularly open.

A useful characterization (in any topological space, hence in $\R$) is:
\[
U\text{ is regularly open }\Longleftrightarrow U=\operatorname{int}(F)\text{ for some closed }F.
\]
Indeed, if $U$ is regularly open, take $F=\overline U$.
Conversely, if $U=\operatorname{int}(F)$ with $F$ closed, then
\[
\overline U\subseteq\overline F=F
\]
and $U\subseteq\operatorname{int}(\overline U)$ since $U$ is open. Also
\[
\operatorname{int}(\overline U)\subseteq\operatorname{int}(F)=U.
\]
Hence $\operatorname{int}(\overline U)=U$.

\textbf{(ii) $(\overline A)^\circ$ is regularly open.}

Set
\[
B=(\overline A)^\circ.
\]
Since $B$ is open, $B\subseteq(\overline B)^\circ$.
Also $B\subseteq\overline A$, so by monotonicity of closure
\[
\overline B\subseteq\overline{\overline A}=\overline A.
\]
Taking interiors,
\[
(\overline B)^\circ\subseteq(\overline A)^\circ=B.
\]
Thus $(\overline B)^\circ=B$, i.e. $B$ is regularly open.
\end{solution}

\section{Question 5}

Let $A\subseteq(X,\mathcal{T})$. Prove
\[
\overline A=A\cup LP(A).
\]

\begin{solution}
First show $A\cup LP(A)\subseteq\overline A$.
Clearly $A\subseteq\overline A$.
If $x\in LP(A)$, every neighborhood $U$ of $x$ meets $A\setminus\{x\}$, hence meets $A$; therefore $x\in\overline A$.
So $A\cup LP(A)\subseteq\overline A$.

Now show $\overline A\subseteq A\cup LP(A)$.
Take $x\in\overline A$.
If $x\in A$, done.
If $x\notin A$, then for any neighborhood $U$ of $x$, because $x\in\overline A$, we have $U\cap A\neq\varnothing$.
Since $x\notin A$, this implies
\[
U\cap(A\setminus\{x\})\neq\varnothing.
\]
Hence $x\in LP(A)$.
So $\overline A\subseteq A\cup LP(A)$.

Combining both inclusions gives
\[
\overline A=A\cup LP(A).
\]
\end{solution}

\end{document}
